\documentclass[11pt]{article}

    \usepackage[breakable]{tcolorbox}
    \usepackage{parskip} % Stop auto-indenting (to mimic markdown behaviour)
    

    % Basic figure setup, for now with no caption control since it's done
    % automatically by Pandoc (which extracts ![](path) syntax from Markdown).
    \usepackage{graphicx}
    % Keep aspect ratio if custom image width or height is specified
    \setkeys{Gin}{keepaspectratio}
    % Maintain compatibility with old templates. Remove in nbconvert 6.0
    \let\Oldincludegraphics\includegraphics
    % Ensure that by default, figures have no caption (until we provide a
    % proper Figure object with a Caption API and a way to capture that
    % in the conversion process - todo).
    \usepackage{caption}
    \DeclareCaptionFormat{nocaption}{}
    \captionsetup{format=nocaption,aboveskip=0pt,belowskip=0pt}

    \usepackage{float}
    \floatplacement{figure}{H} % forces figures to be placed at the correct location
    \usepackage{xcolor} % Allow colors to be defined
    \usepackage{enumerate} % Needed for markdown enumerations to work
    \usepackage{geometry} % Used to adjust the document margins
    \usepackage{amsmath} % Equations
    \usepackage{amssymb} % Equations
    \usepackage{textcomp} % defines textquotesingle
    % Hack from http://tex.stackexchange.com/a/47451/13684:
    \AtBeginDocument{%
        \def\PYZsq{\textquotesingle}% Upright quotes in Pygmentized code
    }
    \usepackage{upquote} % Upright quotes for verbatim code
    \usepackage{eurosym} % defines \euro

    \usepackage{iftex}
    \ifPDFTeX
        \usepackage[T1]{fontenc}
        \IfFileExists{alphabeta.sty}{
              \usepackage{alphabeta}
          }{
              \usepackage[mathletters]{ucs}
              \usepackage[utf8x]{inputenc}
          }
    \else
        \usepackage{fontspec}
        \usepackage{unicode-math}
    \fi

    \usepackage{fancyvrb} % verbatim replacement that allows latex
    \usepackage{grffile} % extends the file name processing of package graphics
                         % to support a larger range
    \makeatletter % fix for old versions of grffile with XeLaTeX
    \@ifpackagelater{grffile}{2019/11/01}
    {
      % Do nothing on new versions
    }
    {
      \def\Gread@@xetex#1{%
        \IfFileExists{"\Gin@base".bb}%
        {\Gread@eps{\Gin@base.bb}}%
        {\Gread@@xetex@aux#1}%
      }
    }
    \makeatother
    \usepackage[Export]{adjustbox} % Used to constrain images to a maximum size
    \adjustboxset{max size={0.9\linewidth}{0.9\paperheight}}

    % The hyperref package gives us a pdf with properly built
    % internal navigation ('pdf bookmarks' for the table of contents,
    % internal cross-reference links, web links for URLs, etc.)
    \usepackage{hyperref}
    % The default LaTeX title has an obnoxious amount of whitespace. By default,
    % titling removes some of it. It also provides customization options.
    \usepackage{titling}
    \usepackage{longtable} % longtable support required by pandoc >1.10
    \usepackage{booktabs}  % table support for pandoc > 1.12.2
    \usepackage{array}     % table support for pandoc >= 2.11.3
    \usepackage{calc}      % table minipage width calculation for pandoc >= 2.11.1
    \usepackage[inline]{enumitem} % IRkernel/repr support (it uses the enumerate* environment)
    \usepackage[normalem]{ulem} % ulem is needed to support strikethroughs (\sout)
                                % normalem makes italics be italics, not underlines
    \usepackage{soul}      % strikethrough (\st) support for pandoc >= 3.0.0
    \usepackage{mathrsfs}
    

    
    % Colors for the hyperref package
    \definecolor{urlcolor}{rgb}{0,.145,.698}
    \definecolor{linkcolor}{rgb}{.71,0.21,0.01}
    \definecolor{citecolor}{rgb}{.12,.54,.11}

    % ANSI colors
    \definecolor{ansi-black}{HTML}{3E424D}
    \definecolor{ansi-black-intense}{HTML}{282C36}
    \definecolor{ansi-red}{HTML}{E75C58}
    \definecolor{ansi-red-intense}{HTML}{B22B31}
    \definecolor{ansi-green}{HTML}{00A250}
    \definecolor{ansi-green-intense}{HTML}{007427}
    \definecolor{ansi-yellow}{HTML}{DDB62B}
    \definecolor{ansi-yellow-intense}{HTML}{B27D12}
    \definecolor{ansi-blue}{HTML}{208FFB}
    \definecolor{ansi-blue-intense}{HTML}{0065CA}
    \definecolor{ansi-magenta}{HTML}{D160C4}
    \definecolor{ansi-magenta-intense}{HTML}{A03196}
    \definecolor{ansi-cyan}{HTML}{60C6C8}
    \definecolor{ansi-cyan-intense}{HTML}{258F8F}
    \definecolor{ansi-white}{HTML}{C5C1B4}
    \definecolor{ansi-white-intense}{HTML}{A1A6B2}
    \definecolor{ansi-default-inverse-fg}{HTML}{FFFFFF}
    \definecolor{ansi-default-inverse-bg}{HTML}{000000}

    % common color for the border for error outputs.
    \definecolor{outerrorbackground}{HTML}{FFDFDF}

    % commands and environments needed by pandoc snippets
    % extracted from the output of `pandoc -s`
    \providecommand{\tightlist}{%
      \setlength{\itemsep}{0pt}\setlength{\parskip}{0pt}}
    \DefineVerbatimEnvironment{Highlighting}{Verbatim}{commandchars=\\\{\}}
    % Add ',fontsize=\small' for more characters per line
    \newenvironment{Shaded}{}{}
    \newcommand{\KeywordTok}[1]{\textcolor[rgb]{0.00,0.44,0.13}{\textbf{{#1}}}}
    \newcommand{\DataTypeTok}[1]{\textcolor[rgb]{0.56,0.13,0.00}{{#1}}}
    \newcommand{\DecValTok}[1]{\textcolor[rgb]{0.25,0.63,0.44}{{#1}}}
    \newcommand{\BaseNTok}[1]{\textcolor[rgb]{0.25,0.63,0.44}{{#1}}}
    \newcommand{\FloatTok}[1]{\textcolor[rgb]{0.25,0.63,0.44}{{#1}}}
    \newcommand{\CharTok}[1]{\textcolor[rgb]{0.25,0.44,0.63}{{#1}}}
    \newcommand{\StringTok}[1]{\textcolor[rgb]{0.25,0.44,0.63}{{#1}}}
    \newcommand{\CommentTok}[1]{\textcolor[rgb]{0.38,0.63,0.69}{\textit{{#1}}}}
    \newcommand{\OtherTok}[1]{\textcolor[rgb]{0.00,0.44,0.13}{{#1}}}
    \newcommand{\AlertTok}[1]{\textcolor[rgb]{1.00,0.00,0.00}{\textbf{{#1}}}}
    \newcommand{\FunctionTok}[1]{\textcolor[rgb]{0.02,0.16,0.49}{{#1}}}
    \newcommand{\RegionMarkerTok}[1]{{#1}}
    \newcommand{\ErrorTok}[1]{\textcolor[rgb]{1.00,0.00,0.00}{\textbf{{#1}}}}
    \newcommand{\NormalTok}[1]{{#1}}

    % Additional commands for more recent versions of Pandoc
    \newcommand{\ConstantTok}[1]{\textcolor[rgb]{0.53,0.00,0.00}{{#1}}}
    \newcommand{\SpecialCharTok}[1]{\textcolor[rgb]{0.25,0.44,0.63}{{#1}}}
    \newcommand{\VerbatimStringTok}[1]{\textcolor[rgb]{0.25,0.44,0.63}{{#1}}}
    \newcommand{\SpecialStringTok}[1]{\textcolor[rgb]{0.73,0.40,0.53}{{#1}}}
    \newcommand{\ImportTok}[1]{{#1}}
    \newcommand{\DocumentationTok}[1]{\textcolor[rgb]{0.73,0.13,0.13}{\textit{{#1}}}}
    \newcommand{\AnnotationTok}[1]{\textcolor[rgb]{0.38,0.63,0.69}{\textbf{\textit{{#1}}}}}
    \newcommand{\CommentVarTok}[1]{\textcolor[rgb]{0.38,0.63,0.69}{\textbf{\textit{{#1}}}}}
    \newcommand{\VariableTok}[1]{\textcolor[rgb]{0.10,0.09,0.49}{{#1}}}
    \newcommand{\ControlFlowTok}[1]{\textcolor[rgb]{0.00,0.44,0.13}{\textbf{{#1}}}}
    \newcommand{\OperatorTok}[1]{\textcolor[rgb]{0.40,0.40,0.40}{{#1}}}
    \newcommand{\BuiltInTok}[1]{{#1}}
    \newcommand{\ExtensionTok}[1]{{#1}}
    \newcommand{\PreprocessorTok}[1]{\textcolor[rgb]{0.74,0.48,0.00}{{#1}}}
    \newcommand{\AttributeTok}[1]{\textcolor[rgb]{0.49,0.56,0.16}{{#1}}}
    \newcommand{\InformationTok}[1]{\textcolor[rgb]{0.38,0.63,0.69}{\textbf{\textit{{#1}}}}}
    \newcommand{\WarningTok}[1]{\textcolor[rgb]{0.38,0.63,0.69}{\textbf{\textit{{#1}}}}}
    \makeatletter
    \newsavebox\pandoc@box
    \newcommand*\pandocbounded[1]{%
      \sbox\pandoc@box{#1}%
      % scaling factors for width and height
      \Gscale@div\@tempa\textheight{\dimexpr\ht\pandoc@box+\dp\pandoc@box\relax}%
      \Gscale@div\@tempb\linewidth{\wd\pandoc@box}%
      % select the smaller of both
      \ifdim\@tempb\p@<\@tempa\p@
        \let\@tempa\@tempb
      \fi
      % scaling accordingly (\@tempa < 1)
      \ifdim\@tempa\p@<\p@
        \scalebox{\@tempa}{\usebox\pandoc@box}%
      % scaling not needed, use as it is
      \else
        \usebox{\pandoc@box}%
      \fi
    }
    \makeatother

    % Define a nice break command that doesn't care if a line doesn't already
    % exist.
    \def\br{\hspace*{\fill} \\* }
    % Math Jax compatibility definitions
    \def\gt{>}
    \def\lt{<}
    \let\Oldtex\TeX
    \let\Oldlatex\LaTeX
    \renewcommand{\TeX}{\textrm{\Oldtex}}
    \renewcommand{\LaTeX}{\textrm{\Oldlatex}}
    % Document parameters
    % Document title
    \title{building\_silver\_layer\_v2}
    
    
    
    
    
    
    
% Pygments definitions
\makeatletter
\def\PY@reset{\let\PY@it=\relax \let\PY@bf=\relax%
    \let\PY@ul=\relax \let\PY@tc=\relax%
    \let\PY@bc=\relax \let\PY@ff=\relax}
\def\PY@tok#1{\csname PY@tok@#1\endcsname}
\def\PY@toks#1+{\ifx\relax#1\empty\else%
    \PY@tok{#1}\expandafter\PY@toks\fi}
\def\PY@do#1{\PY@bc{\PY@tc{\PY@ul{%
    \PY@it{\PY@bf{\PY@ff{#1}}}}}}}
\def\PY#1#2{\PY@reset\PY@toks#1+\relax+\PY@do{#2}}

\@namedef{PY@tok@w}{\def\PY@tc##1{\textcolor[rgb]{0.73,0.73,0.73}{##1}}}
\@namedef{PY@tok@c}{\let\PY@it=\textit\def\PY@tc##1{\textcolor[rgb]{0.24,0.48,0.48}{##1}}}
\@namedef{PY@tok@cp}{\def\PY@tc##1{\textcolor[rgb]{0.61,0.40,0.00}{##1}}}
\@namedef{PY@tok@k}{\let\PY@bf=\textbf\def\PY@tc##1{\textcolor[rgb]{0.00,0.50,0.00}{##1}}}
\@namedef{PY@tok@kp}{\def\PY@tc##1{\textcolor[rgb]{0.00,0.50,0.00}{##1}}}
\@namedef{PY@tok@kt}{\def\PY@tc##1{\textcolor[rgb]{0.69,0.00,0.25}{##1}}}
\@namedef{PY@tok@o}{\def\PY@tc##1{\textcolor[rgb]{0.40,0.40,0.40}{##1}}}
\@namedef{PY@tok@ow}{\let\PY@bf=\textbf\def\PY@tc##1{\textcolor[rgb]{0.67,0.13,1.00}{##1}}}
\@namedef{PY@tok@nb}{\def\PY@tc##1{\textcolor[rgb]{0.00,0.50,0.00}{##1}}}
\@namedef{PY@tok@nf}{\def\PY@tc##1{\textcolor[rgb]{0.00,0.00,1.00}{##1}}}
\@namedef{PY@tok@nc}{\let\PY@bf=\textbf\def\PY@tc##1{\textcolor[rgb]{0.00,0.00,1.00}{##1}}}
\@namedef{PY@tok@nn}{\let\PY@bf=\textbf\def\PY@tc##1{\textcolor[rgb]{0.00,0.00,1.00}{##1}}}
\@namedef{PY@tok@ne}{\let\PY@bf=\textbf\def\PY@tc##1{\textcolor[rgb]{0.80,0.25,0.22}{##1}}}
\@namedef{PY@tok@nv}{\def\PY@tc##1{\textcolor[rgb]{0.10,0.09,0.49}{##1}}}
\@namedef{PY@tok@no}{\def\PY@tc##1{\textcolor[rgb]{0.53,0.00,0.00}{##1}}}
\@namedef{PY@tok@nl}{\def\PY@tc##1{\textcolor[rgb]{0.46,0.46,0.00}{##1}}}
\@namedef{PY@tok@ni}{\let\PY@bf=\textbf\def\PY@tc##1{\textcolor[rgb]{0.44,0.44,0.44}{##1}}}
\@namedef{PY@tok@na}{\def\PY@tc##1{\textcolor[rgb]{0.41,0.47,0.13}{##1}}}
\@namedef{PY@tok@nt}{\let\PY@bf=\textbf\def\PY@tc##1{\textcolor[rgb]{0.00,0.50,0.00}{##1}}}
\@namedef{PY@tok@nd}{\def\PY@tc##1{\textcolor[rgb]{0.67,0.13,1.00}{##1}}}
\@namedef{PY@tok@s}{\def\PY@tc##1{\textcolor[rgb]{0.73,0.13,0.13}{##1}}}
\@namedef{PY@tok@sd}{\let\PY@it=\textit\def\PY@tc##1{\textcolor[rgb]{0.73,0.13,0.13}{##1}}}
\@namedef{PY@tok@si}{\let\PY@bf=\textbf\def\PY@tc##1{\textcolor[rgb]{0.64,0.35,0.47}{##1}}}
\@namedef{PY@tok@se}{\let\PY@bf=\textbf\def\PY@tc##1{\textcolor[rgb]{0.67,0.36,0.12}{##1}}}
\@namedef{PY@tok@sr}{\def\PY@tc##1{\textcolor[rgb]{0.64,0.35,0.47}{##1}}}
\@namedef{PY@tok@ss}{\def\PY@tc##1{\textcolor[rgb]{0.10,0.09,0.49}{##1}}}
\@namedef{PY@tok@sx}{\def\PY@tc##1{\textcolor[rgb]{0.00,0.50,0.00}{##1}}}
\@namedef{PY@tok@m}{\def\PY@tc##1{\textcolor[rgb]{0.40,0.40,0.40}{##1}}}
\@namedef{PY@tok@gh}{\let\PY@bf=\textbf\def\PY@tc##1{\textcolor[rgb]{0.00,0.00,0.50}{##1}}}
\@namedef{PY@tok@gu}{\let\PY@bf=\textbf\def\PY@tc##1{\textcolor[rgb]{0.50,0.00,0.50}{##1}}}
\@namedef{PY@tok@gd}{\def\PY@tc##1{\textcolor[rgb]{0.63,0.00,0.00}{##1}}}
\@namedef{PY@tok@gi}{\def\PY@tc##1{\textcolor[rgb]{0.00,0.52,0.00}{##1}}}
\@namedef{PY@tok@gr}{\def\PY@tc##1{\textcolor[rgb]{0.89,0.00,0.00}{##1}}}
\@namedef{PY@tok@ge}{\let\PY@it=\textit}
\@namedef{PY@tok@gs}{\let\PY@bf=\textbf}
\@namedef{PY@tok@ges}{\let\PY@bf=\textbf\let\PY@it=\textit}
\@namedef{PY@tok@gp}{\let\PY@bf=\textbf\def\PY@tc##1{\textcolor[rgb]{0.00,0.00,0.50}{##1}}}
\@namedef{PY@tok@go}{\def\PY@tc##1{\textcolor[rgb]{0.44,0.44,0.44}{##1}}}
\@namedef{PY@tok@gt}{\def\PY@tc##1{\textcolor[rgb]{0.00,0.27,0.87}{##1}}}
\@namedef{PY@tok@err}{\def\PY@bc##1{{\setlength{\fboxsep}{\string -\fboxrule}\fcolorbox[rgb]{1.00,0.00,0.00}{1,1,1}{\strut ##1}}}}
\@namedef{PY@tok@kc}{\let\PY@bf=\textbf\def\PY@tc##1{\textcolor[rgb]{0.00,0.50,0.00}{##1}}}
\@namedef{PY@tok@kd}{\let\PY@bf=\textbf\def\PY@tc##1{\textcolor[rgb]{0.00,0.50,0.00}{##1}}}
\@namedef{PY@tok@kn}{\let\PY@bf=\textbf\def\PY@tc##1{\textcolor[rgb]{0.00,0.50,0.00}{##1}}}
\@namedef{PY@tok@kr}{\let\PY@bf=\textbf\def\PY@tc##1{\textcolor[rgb]{0.00,0.50,0.00}{##1}}}
\@namedef{PY@tok@bp}{\def\PY@tc##1{\textcolor[rgb]{0.00,0.50,0.00}{##1}}}
\@namedef{PY@tok@fm}{\def\PY@tc##1{\textcolor[rgb]{0.00,0.00,1.00}{##1}}}
\@namedef{PY@tok@vc}{\def\PY@tc##1{\textcolor[rgb]{0.10,0.09,0.49}{##1}}}
\@namedef{PY@tok@vg}{\def\PY@tc##1{\textcolor[rgb]{0.10,0.09,0.49}{##1}}}
\@namedef{PY@tok@vi}{\def\PY@tc##1{\textcolor[rgb]{0.10,0.09,0.49}{##1}}}
\@namedef{PY@tok@vm}{\def\PY@tc##1{\textcolor[rgb]{0.10,0.09,0.49}{##1}}}
\@namedef{PY@tok@sa}{\def\PY@tc##1{\textcolor[rgb]{0.73,0.13,0.13}{##1}}}
\@namedef{PY@tok@sb}{\def\PY@tc##1{\textcolor[rgb]{0.73,0.13,0.13}{##1}}}
\@namedef{PY@tok@sc}{\def\PY@tc##1{\textcolor[rgb]{0.73,0.13,0.13}{##1}}}
\@namedef{PY@tok@dl}{\def\PY@tc##1{\textcolor[rgb]{0.73,0.13,0.13}{##1}}}
\@namedef{PY@tok@s2}{\def\PY@tc##1{\textcolor[rgb]{0.73,0.13,0.13}{##1}}}
\@namedef{PY@tok@sh}{\def\PY@tc##1{\textcolor[rgb]{0.73,0.13,0.13}{##1}}}
\@namedef{PY@tok@s1}{\def\PY@tc##1{\textcolor[rgb]{0.73,0.13,0.13}{##1}}}
\@namedef{PY@tok@mb}{\def\PY@tc##1{\textcolor[rgb]{0.40,0.40,0.40}{##1}}}
\@namedef{PY@tok@mf}{\def\PY@tc##1{\textcolor[rgb]{0.40,0.40,0.40}{##1}}}
\@namedef{PY@tok@mh}{\def\PY@tc##1{\textcolor[rgb]{0.40,0.40,0.40}{##1}}}
\@namedef{PY@tok@mi}{\def\PY@tc##1{\textcolor[rgb]{0.40,0.40,0.40}{##1}}}
\@namedef{PY@tok@il}{\def\PY@tc##1{\textcolor[rgb]{0.40,0.40,0.40}{##1}}}
\@namedef{PY@tok@mo}{\def\PY@tc##1{\textcolor[rgb]{0.40,0.40,0.40}{##1}}}
\@namedef{PY@tok@ch}{\let\PY@it=\textit\def\PY@tc##1{\textcolor[rgb]{0.24,0.48,0.48}{##1}}}
\@namedef{PY@tok@cm}{\let\PY@it=\textit\def\PY@tc##1{\textcolor[rgb]{0.24,0.48,0.48}{##1}}}
\@namedef{PY@tok@cpf}{\let\PY@it=\textit\def\PY@tc##1{\textcolor[rgb]{0.24,0.48,0.48}{##1}}}
\@namedef{PY@tok@c1}{\let\PY@it=\textit\def\PY@tc##1{\textcolor[rgb]{0.24,0.48,0.48}{##1}}}
\@namedef{PY@tok@cs}{\let\PY@it=\textit\def\PY@tc##1{\textcolor[rgb]{0.24,0.48,0.48}{##1}}}

\def\PYZbs{\char`\\}
\def\PYZus{\char`\_}
\def\PYZob{\char`\{}
\def\PYZcb{\char`\}}
\def\PYZca{\char`\^}
\def\PYZam{\char`\&}
\def\PYZlt{\char`\<}
\def\PYZgt{\char`\>}
\def\PYZsh{\char`\#}
\def\PYZpc{\char`\%}
\def\PYZdl{\char`\$}
\def\PYZhy{\char`\-}
\def\PYZsq{\char`\'}
\def\PYZdq{\char`\"}
\def\PYZti{\char`\~}
% for compatibility with earlier versions
\def\PYZat{@}
\def\PYZlb{[}
\def\PYZrb{]}
\makeatother


    % For linebreaks inside Verbatim environment from package fancyvrb.
    \makeatletter
        \newbox\Wrappedcontinuationbox
        \newbox\Wrappedvisiblespacebox
        \newcommand*\Wrappedvisiblespace {\textcolor{red}{\textvisiblespace}}
        \newcommand*\Wrappedcontinuationsymbol {\textcolor{red}{\llap{\tiny$\m@th\hookrightarrow$}}}
        \newcommand*\Wrappedcontinuationindent {3ex }
        \newcommand*\Wrappedafterbreak {\kern\Wrappedcontinuationindent\copy\Wrappedcontinuationbox}
        % Take advantage of the already applied Pygments mark-up to insert
        % potential linebreaks for TeX processing.
        %        {, <, #, %, $, ' and ": go to next line.
        %        _, }, ^, &, >, - and ~: stay at end of broken line.
        % Use of \textquotesingle for straight quote.
        \newcommand*\Wrappedbreaksatspecials {%
            \def\PYGZus{\discretionary{\char`\_}{\Wrappedafterbreak}{\char`\_}}%
            \def\PYGZob{\discretionary{}{\Wrappedafterbreak\char`\{}{\char`\{}}%
            \def\PYGZcb{\discretionary{\char`\}}{\Wrappedafterbreak}{\char`\}}}%
            \def\PYGZca{\discretionary{\char`\^}{\Wrappedafterbreak}{\char`\^}}%
            \def\PYGZam{\discretionary{\char`\&}{\Wrappedafterbreak}{\char`\&}}%
            \def\PYGZlt{\discretionary{}{\Wrappedafterbreak\char`\<}{\char`\<}}%
            \def\PYGZgt{\discretionary{\char`\>}{\Wrappedafterbreak}{\char`\>}}%
            \def\PYGZsh{\discretionary{}{\Wrappedafterbreak\char`\#}{\char`\#}}%
            \def\PYGZpc{\discretionary{}{\Wrappedafterbreak\char`\%}{\char`\%}}%
            \def\PYGZdl{\discretionary{}{\Wrappedafterbreak\char`\$}{\char`\$}}%
            \def\PYGZhy{\discretionary{\char`\-}{\Wrappedafterbreak}{\char`\-}}%
            \def\PYGZsq{\discretionary{}{\Wrappedafterbreak\textquotesingle}{\textquotesingle}}%
            \def\PYGZdq{\discretionary{}{\Wrappedafterbreak\char`\"}{\char`\"}}%
            \def\PYGZti{\discretionary{\char`\~}{\Wrappedafterbreak}{\char`\~}}%
        }
        % Some characters . , ; ? ! / are not pygmentized.
        % This macro makes them "active" and they will insert potential linebreaks
        \newcommand*\Wrappedbreaksatpunct {%
            \lccode`\~`\.\lowercase{\def~}{\discretionary{\hbox{\char`\.}}{\Wrappedafterbreak}{\hbox{\char`\.}}}%
            \lccode`\~`\,\lowercase{\def~}{\discretionary{\hbox{\char`\,}}{\Wrappedafterbreak}{\hbox{\char`\,}}}%
            \lccode`\~`\;\lowercase{\def~}{\discretionary{\hbox{\char`\;}}{\Wrappedafterbreak}{\hbox{\char`\;}}}%
            \lccode`\~`\:\lowercase{\def~}{\discretionary{\hbox{\char`\:}}{\Wrappedafterbreak}{\hbox{\char`\:}}}%
            \lccode`\~`\?\lowercase{\def~}{\discretionary{\hbox{\char`\?}}{\Wrappedafterbreak}{\hbox{\char`\?}}}%
            \lccode`\~`\!\lowercase{\def~}{\discretionary{\hbox{\char`\!}}{\Wrappedafterbreak}{\hbox{\char`\!}}}%
            \lccode`\~`\/\lowercase{\def~}{\discretionary{\hbox{\char`\/}}{\Wrappedafterbreak}{\hbox{\char`\/}}}%
            \catcode`\.\active
            \catcode`\,\active
            \catcode`\;\active
            \catcode`\:\active
            \catcode`\?\active
            \catcode`\!\active
            \catcode`\/\active
            \lccode`\~`\~
        }
    \makeatother

    \let\OriginalVerbatim=\Verbatim
    \makeatletter
    \renewcommand{\Verbatim}[1][1]{%
        %\parskip\z@skip
        \sbox\Wrappedcontinuationbox {\Wrappedcontinuationsymbol}%
        \sbox\Wrappedvisiblespacebox {\FV@SetupFont\Wrappedvisiblespace}%
        \def\FancyVerbFormatLine ##1{\hsize\linewidth
            \vtop{\raggedright\hyphenpenalty\z@\exhyphenpenalty\z@
                \doublehyphendemerits\z@\finalhyphendemerits\z@
                \strut ##1\strut}%
        }%
        % If the linebreak is at a space, the latter will be displayed as visible
        % space at end of first line, and a continuation symbol starts next line.
        % Stretch/shrink are however usually zero for typewriter font.
        \def\FV@Space {%
            \nobreak\hskip\z@ plus\fontdimen3\font minus\fontdimen4\font
            \discretionary{\copy\Wrappedvisiblespacebox}{\Wrappedafterbreak}
            {\kern\fontdimen2\font}%
        }%

        % Allow breaks at special characters using \PYG... macros.
        \Wrappedbreaksatspecials
        % Breaks at punctuation characters . , ; ? ! and / need catcode=\active
        \OriginalVerbatim[#1,codes*=\Wrappedbreaksatpunct]%
    }
    \makeatother

    % Exact colors from NB
    \definecolor{incolor}{HTML}{303F9F}
    \definecolor{outcolor}{HTML}{D84315}
    \definecolor{cellborder}{HTML}{CFCFCF}
    \definecolor{cellbackground}{HTML}{F7F7F7}

    % prompt
    \makeatletter
    \newcommand{\boxspacing}{\kern\kvtcb@left@rule\kern\kvtcb@boxsep}
    \makeatother
    \newcommand{\prompt}[4]{
        {\ttfamily\llap{{\color{#2}[#3]:\hspace{3pt}#4}}\vspace{-\baselineskip}}
    }
    

    
    % Prevent overflowing lines due to hard-to-break entities
    \sloppy
    % Setup hyperref package
    \hypersetup{
      breaklinks=true,  % so long urls are correctly broken across lines
      colorlinks=true,
      urlcolor=urlcolor,
      linkcolor=linkcolor,
      citecolor=citecolor,
      }
    % Slightly bigger margins than the latex defaults
    
    \geometry{verbose,tmargin=1in,bmargin=1in,lmargin=1in,rmargin=1in}
    
    

\begin{document}
    
    \maketitle
    
    

    
    \section{Kolumbus Scheduled and Real-Time Bus Service: Analysis
Report}\label{kolumbus-scheduled-and-real-time-bus-service-analysis-report}

\subsection{Problem and aim}\label{problem-and-aim}

I am investigating how the bus service in the published GTFS timetable
compares with the stop-call information recorded in the SIRI Estimated
Timetable (ET) real-time feed. The problem is that these sources
describe service in different ways: GTFS says what is scheduled, while
the real-time feed records changing predictions and, when available,
what actually happened. The real-time data also contains repeated
snapshots, so its raw row counts cannot be treated as counts of bus
visits.

My aim is to build a reliable record of scheduled and real-time bus stop
events, then assess which scheduled events are confirmed, cancelled,
unconfirmed, or not represented in the feed. I want to understand the
coverage and quality of the data before drawing conclusions about bus
operations or delays. In particular, a missing real-time record is not
by itself proof that a bus did not run.

\subsection{What I am doing}\label{what-i-am-doing}

I first inspect the Parquet files and their schemas. I then use the
static GTFS relationships between routes, trips, stops, stop times, and
service calendars to reconstruct the bus stop events expected on each
service date. I review service alerts and check that their identifiers
can be related to GTFS records.

Next, I profile the real-time \texttt{stop\_visits} data. Because one
journey-stop event may appear in multiple snapshots as predictions
change, I consolidate those records into one visit, using the latest
snapshot while retaining actual arrival or departure times that may have
appeared earlier. I validate the records against GTFS, limit the
analysis to bus services and known stops, and classify each visit by its
available evidence.

Finally, I match scheduled bus stop events with real-time visits using
the service date, trip, stop, and occurrence of that stop in the trip. I
use the matched records to summarize scheduled and observed activity by
stop, route, and date, and to calculate delay where an actual time is
available.

\subsection{How I interpret the
results}\label{how-i-interpret-the-results}

\begin{itemize}
\tightlist
\item
  \textbf{Expected} means the GTFS timetable and service calendar
  indicate that the stop event was scheduled.
\item
  \textbf{Observed} means a real-time record contains an actual arrival
  or departure time after snapshots have been consolidated.
\item
  \textbf{Cancelled} means the real-time feed marks the journey or stop
  call as cancelled.
\item
  \textbf{Unconfirmed} means a real-time record exists but contains no
  actual arrival or departure time.
\item
  \textbf{Not in feed} means no matching real-time record was found.
  This may indicate incomplete feed coverage and is not automatically a
  missed service.
\end{itemize}

The conclusions are limited by the dates available, the accuracy of the
source identifiers and timetable, and the completeness of the real-time
feed. I calculate delay only when an actual time can be compared with
its aimed timetable time. I keep timestamps in UTC and derive local-time
fields for interpretation, taking daylight-saving changes into account.

\subsection{Analysis structure}\label{analysis-structure}

\begin{longtable}[]{@{}
  >{\raggedright\arraybackslash}p{(\linewidth - 2\tabcolsep) * \real{0.5000}}
  >{\raggedright\arraybackslash}p{(\linewidth - 2\tabcolsep) * \real{0.5000}}@{}}
\toprule\noalign{}
\begin{minipage}[b]{\linewidth}\raggedright
Part
\end{minipage} & \begin{minipage}[b]{\linewidth}\raggedright
Analysis
\end{minipage} \\
\midrule\noalign{}
\endhead
\bottomrule\noalign{}
\endlastfoot
0?1 & Configure the analysis and inspect the source files \\
2 & Examine static GTFS routes, trips, stops, and timetable
relationships \\
3 & Review service alerts and check their links to GTFS identifiers \\
4 & Profile the real-time feed, its repeated snapshots, and data
quality \\
5 & Reconstruct scheduled bus stop events from the timetable and service
calendars \\
6 & Clean, validate, and classify real-time stop visits \\
7 & Compare scheduled events with real-time outcomes by stop and
route \\
8 & Save the cleaned and comparison tables for further analysis \\
\end{longtable}

\textbf{Important:} I consolidate repeated snapshots before counting
visits or comparing observed events with the timetable.

    \subsection{0. Analysis setup}\label{analysis-setup}

I set the data location, output folder, local timezone, and service-date
selection used throughout the notebook.

    \subsection{1. Source-file overview}\label{source-file-overview}

I begin by checking the available Parquet files, their row and column
counts, and the snapshot dates in the static GTFS tables. This
establishes what data is available before I build the timetable or
compare it with real-time records.

    
    \begin{Verbatim}[commandchars=\\\{\}]
                   file     rows  cols  row\_groups
0           gtfs\_agency        3     6           1
1           gtfs\_routes      318     7           1
2            gtfs\_trips    13833     9           1
3       gtfs\_stop\_times   418837    11           1
4            gtfs\_stops     7676    12           1
5         gtfs\_calendar      223    12           1
6   gtfs\_calendar\_dates     2722     5           1
7           gtfs\_shapes   614593     7           1
8        service\_alerts       53    17           1
9           stop\_visits  9989935    26         100
10    vehicle\_positions   300316    13           4
    \end{Verbatim}

    
    \begin{Verbatim}[commandchars=\\\{\}]
gtfs\_agency          agency\_id, agency\_name, agency\_url, agency\_timezone,
snapshot\_time, source\_file
gtfs\_routes          agency\_id, route\_id, route\_short\_name, route\_long\_name,
route\_type, snapshot\_time, source\_file
gtfs\_trips           route\_id, trip\_id, service\_id, trip\_headsign, direction\_id,
shape\_id, boarding\_type, snapshot\_time, source\_file
gtfs\_stop\_times      trip\_id, stop\_id, arrival\_time, departure\_time,
stop\_sequence, stop\_headsign, pickup\_type, drop\_off\_type, shape\_dist\_traveled,
snapshot\_time, source\_file
gtfs\_stops           stop\_id, stop\_name, stop\_lat, stop\_lon, stop\_desc,
location\_type, parent\_station, wheelchair\_boarding, vehicle\_type, platform\_code,
snapshot\_time, source\_file
gtfs\_calendar        service\_id, monday, tuesday, wednesday, thursday, friday,
saturday, sunday, start\_date, end\_date, snapshot\_time, source\_file
gtfs\_calendar\_dates  service\_id, date, exception\_type, snapshot\_time,
source\_file
gtfs\_shapes          shape\_id, shape\_pt\_sequence, shape\_pt\_lat, shape\_pt\_lon,
shape\_dist\_traveled, snapshot\_time, source\_file
service\_alerts       snapshot\_time, situation\_number, creation\_time,
version\_time, progress, severity, report\_type, planned, summary, valid\_start,
valid\_end, affect\_type, line\_ref, journey\_ref, operator\_ref,
origin\_aimed\_departure, source\_file
stop\_visits          snapshot\_time, recorded\_at, journey\_ref, data\_frame\_ref,
line\_ref, direction\_ref, vehicle\_mode, operator\_ref, monitored,
journey\_cancelled, call\_type, stop\_order, stop\_ref, stop\_name, destination,
aimed\_arrival, expected\_arrival, actual\_arrival, aimed\_departure,
expected\_departure, actual\_departure, arrival\_status, departure\_status,
call\_cancelled, prediction\_inaccurate, source\_file
vehicle\_positions    snapshot\_time, vehicle\_id, last\_updated, line\_ref,
line\_name, public\_code, latitude, longitude, bearing, speed, delay,
occupancy\_status, source\_file
    \end{Verbatim}

    \textbf{Data considerations.} The GTFS fields are stored as strings,
including dates, sequence values, and flags. I convert these fields to
appropriate types before using them in calculations or comparisons. The
\texttt{snapshot\_time} and \texttt{source\_file} fields were added by
the data pipeline rather than being standard GTFS fields. I check the
snapshot metadata to confirm that the static tables represent a
consistent timetable snapshot and are not duplicated across multiple
extracts.

    \begin{Verbatim}[commandchars=\\\{\}]
gtfs\_agency          distinct snapshots: 1  (2026-10-24 20:00:12+00:00 →
2026-10-24 20:00:12+00:00)
gtfs\_routes          distinct snapshots: 1  (2026-10-24 20:00:12+00:00 →
2026-10-24 20:00:12+00:00)
gtfs\_trips           distinct snapshots: 1  (2026-10-24 20:00:12+00:00 →
2026-10-24 20:00:12+00:00)
gtfs\_stop\_times      distinct snapshots: 1  (2026-10-24 20:00:12+00:00 →
2026-10-24 20:00:12+00:00)
gtfs\_stops           distinct snapshots: 1  (2026-10-24 20:00:12+00:00 →
2026-10-24 20:00:12+00:00)
gtfs\_calendar        distinct snapshots: 1  (2026-10-24 20:00:12+00:00 →
2026-10-24 20:00:12+00:00)
gtfs\_calendar\_dates  distinct snapshots: 1  (2026-10-24 20:00:12+00:00 →
2026-10-24 20:00:12+00:00)
    \end{Verbatim}

    \subsection{2. Static GTFS: the scheduled
service}\label{static-gtfs-the-scheduled-service}

I use the static GTFS tables to describe the service the timetable says
should operate. The route, trip, stop, stop-time, and calendar tables
together provide the basis for identifying scheduled bus events.

    
    \begin{Verbatim}[commandchars=\\\{\}]
          agency\_id agency\_name agency\_timezone
0  KOL:Authority:26   Norled AS     Europe/Oslo
1  KOL:Authority:30  Fjord1 ASA     Europe/Oslo
2   KOL:Authority:8    Kolumbus     Europe/Oslo
    \end{Verbatim}

    
    \begin{Verbatim}[commandchars=\\\{\}]
318 routes | 13,833 trips | 418,837 stop\_times | 7,676 stops
    \end{Verbatim}

    \subsubsection{\texorpdfstring{2.1 Interpreting
\texttt{route\_type}}{2.1 Interpreting route\_type}}\label{interpreting-route_type}

The route types in these files use the extended GTFS codes rather than
only the basic 0?7 values. I use the following common interpretations to
distinguish bus services from ferry services. If a route-type label
affects a reported conclusion, I verify it against the applicable
official GTFS reference.

            \begin{tcolorbox}[breakable, size=fbox, boxrule=.5pt, pad at break*=1mm, opacityfill=0]
\begin{Verbatim}[commandchars=\\\{\}]
  route\_type             route\_type\_name  n\_routes
5        712                  School bus       179
3        704                   Local bus       105
4        705                   Night bus        13
2       1014  High-speed passenger ferry        12
0       1004             Local car ferry         8
1       1008       Local passenger ferry         1
\end{Verbatim}
\end{tcolorbox}
        
    \subsubsection{2.2 Checking the school-bus
classification}\label{checking-the-school-bus-classification}

I check whether route type \texttt{712} behaves like a school-bus
category by comparing the stops served by routes of type \texttt{712}
with those served by local bus routes of type \texttt{704}. This is a
more systematic check than relying on individual route numbers. I expect
the service patterns to differ, although major terminals may be shared.
The comparison provides evidence for interpreting the category; it does
not, by itself, prove the intended meaning of the code.

    \begin{Verbatim}[commandchars=\\\{\}]
stops served by 704 (local bus):  3,821
stops served by 712 (school bus): 3,074
stops served by BOTH:             2,202
stops served ONLY by 712:         872
    \end{Verbatim}

    
    \begin{Verbatim}[commandchars=\\\{\}]
                  stop\_id stop\_name        served\_by
1097       NSR:Quay:44969  UiS nord  [704, 705, 712]
1098       NSR:Quay:44970  UiS nord       [704, 705]
1108       NSR:Quay:44984   UiS øst  [704, 705, 712]
1110       NSR:Quay:44986   UiS øst       [704, 705]
4930  NSR:StopPlace:25943  UiS nord               []
4937  NSR:StopPlace:25953   UiS øst               []
    \end{Verbatim}

    
    \subsubsection{2.3 Linking a shape to trips, routes, and
stops}\label{linking-a-shape-to-trips-routes-and-stops}

A \texttt{shape\_id} describes a path, but it does not alone identify
the timetable's ordered stop calls. I trace the relationships from
\texttt{shape\_id} to trips, from trips to routes, and from each trip to
its ordered \texttt{stop\_times}; I then use the stops table to retrieve
names and coordinates. The function below makes this lookup repeatable
for any shape I want to examine.

    \begin{Verbatim}[commandchars=\\\{\}]
248 trips use this shape, on route(s):
    \end{Verbatim}

    
    \begin{Verbatim}[commandchars=\\\{\}]
          route\_id route\_short\_name route\_long\_name route\_type route\_type\_name
2  KOL:Line:8\_1001                1                        704       Local bus
    \end{Verbatim}

    
    \begin{Verbatim}[commandchars=\\\{\}]
Stops for example trip KOL:ServiceJourney:1001\_251008118370980\_2001:
    \end{Verbatim}

            \begin{tcolorbox}[breakable, size=fbox, boxrule=.5pt, pad at break*=1mm, opacityfill=0]
\begin{Verbatim}[commandchars=\\\{\}]
    stop\_sequence          stop\_id                   stop\_name arrival\_time
departure\_time
0               1  NSR:Quay:103376  Sandnes bussterminal Ruten     05:30:00
05:30:00
1               2  NSR:Quay:109015               Olav Vs plass     05:32:00
05:32:00
2               3  NSR:Quay:111020               Jernbaneveien     05:33:00
05:33:00
3               4   NSR:Quay:47663                  Strandgata     05:35:00
05:35:00
4               5  NSR:Quay:109013             Strandgata nord     05:37:00
05:37:00
5               6   NSR:Quay:49004               Gamle Postvei     05:39:00
05:39:00
6               7   NSR:Quay:45344                  Seljeveien     05:40:00
05:40:00
7               8   NSR:Quay:47836                     Kvadrat     05:41:00
05:41:00
8               9   NSR:Quay:48301                Porsmyrveien     05:42:00
05:42:00
9              10   NSR:Quay:46365         Jakob Askelands vei     05:43:00
05:43:00
10             11   NSR:Quay:47231                Løwenstrasse     05:44:00
05:44:00
11             12   NSR:Quay:48774                   2020 park     05:45:00
05:45:00
12             13   NSR:Quay:48992              Nådlandsbråtet     05:46:00
05:46:00
13             14   NSR:Quay:49080              Gausel stasjon     05:47:00
05:47:00
14             15   NSR:Quay:49085                 Gauselhagen     05:49:00
05:49:00
15             16   NSR:Quay:49078              Gausel sentrum     05:50:00
05:50:00
16             17  NSR:Quay:108991                       Jåttå     05:52:00
05:52:00
17             18  NSR:Quay:108996               Stadionparken     05:53:00
05:53:00
18             19  NSR:Quay:108998                  Skråtårnet     05:54:00
05:54:00
19             20  NSR:Quay:108999               Hinnasvingene     05:56:00
05:56:00
20             21   NSR:Quay:47141                Lyngnesveien     05:58:00
05:58:00
21             22   NSR:Quay:47331                     Mariero     05:59:00
05:59:00
22             23   NSR:Quay:48151                    Eikeberg     06:00:00
06:00:00
23             24   NSR:Quay:45392                    Sjøhagen     06:01:00
06:01:00
24             25   NSR:Quay:47870             Kvaleberg skole     06:02:00
06:02:00
25             26   NSR:Quay:45727                    Hillevåg     06:03:00
06:03:00
26             27   NSR:Quay:45732           Hillevågstunnelen     06:04:00
06:04:00
27             28   NSR:Quay:47270                  Strømsbrua     06:05:00
06:05:00
28             29   NSR:Quay:46034               Vålandsbakken     06:06:00
06:06:00
29             30   NSR:Quay:46919                 Statens hus     06:07:00
06:07:00
30             31   NSR:Quay:47005     Stavanger politistasjon     06:08:00
06:08:00
31             32   NSR:Quay:46974               Jernbaneveien     06:09:00
06:09:00
\end{Verbatim}
\end{tcolorbox}
        
    \subsubsection{2.4 Checking static-table
integrity}\label{checking-static-table-integrity}

Before using the timetable to create expected stop events, I check for
duplicate identifiers and broken references between trips, stop times,
and stops. A non-zero problem count indicates records that need
investigation because they may affect joins or event counts.

    \begin{Verbatim}[commandchars=\\\{\}]
OK   duplicate trip\_id in trips: 0 of 13,833
OK   duplicate stop\_id in stops: 0 of 7,676
OK   stop\_times with trip\_id not in trips: 0 of 418,837
OK   stop\_times with stop\_id not in stops: 0 of 418,837
OK   trips with route\_id not in routes: 0 of 13,833
OK   trips with service\_id in neither calendar nor calendar\_dates: 0 of 13,833
OK   stops with missing coordinates: 0 of 7,676
calendar service dates: 20261016 → 20270403
exception types in calendar\_dates: \{'2': 1942, '1': 780\} (1 = service added, 2 =
service removed)
    \end{Verbatim}

    \subsection{3. Service alerts}\label{service-alerts}

I examine the service-alert records as additional context for
interpreting service disruptions. The relevant category field in this
dataset is \texttt{affect\_type}, with values such as \texttt{line} and
\texttt{journey}. Alerts can be repeated across snapshots, so I retain
the latest record for each alert reference before summarizing or
matching alerts to timetable entities.

    \begin{Verbatim}[commandchars=\\\{\}]
53 alert rows, 10 distinct situations
affect\_type
line       44
journey     9

severity
normal      41
noImpact     6
severe       6

summary
Mulige forsinkelser pga. vanskelige værforhold    41
Forsinket på grunn av mye trafikk                  4
Trafikkmelding                                     3
Redusert kapasitet                                 3
Forsinket på grunn av tekniske problemer           1
Tur avsluttes før endeholdeplass                   1

after de-duplicating snapshots: 11 rows
affect\_type
journey    6
line       5

    \end{Verbatim}

    
    \begin{Verbatim}[commandchars=\\\{\}]
               snapshot\_time                                   summary  severity          line\_ref                                   journey\_ref                valid\_start                  valid\_end
9  2026-10-24 10:59:35+00:00                        Redusert kapasitet  noImpact  KOL:Line:26\_5958                                           NaN  2026-10-24T13:00:00+02:00  2026-10-24T13:12:00+02:00
16 2026-10-24 11:51:39+00:00  Forsinket på grunn av tekniske problemer    severe   KOL:Line:8\_1006  KOL:ServiceJourney:1006\_261005064436180\_3024  2026-10-24T00:00:00+02:00  2026-10-24T15:57:00+02:00
28 2026-10-24 15:57:00+00:00         Forsinket på grunn av mye trafikk    severe   KOL:Line:8\_2010  KOL:ServiceJourney:2010\_260923062817478\_4014  2026-10-24T00:00:00+02:00  2026-10-24T18:06:00+02:00
32 2026-10-24 16:17:01+00:00         Forsinket på grunn av mye trafikk    severe   KOL:Line:8\_6100  KOL:ServiceJourney:6100\_250903083239593\_4017  2026-10-24T00:00:00+02:00  2026-10-24T19:28:00+02:00
43 2026-10-24 19:39:18+00:00                        Redusert kapasitet  noImpact  KOL:Line:26\_5958                                           NaN  2026-10-24T21:40:00+02:00  2026-10-24T21:53:00+02:00
50 2026-10-24 23:24:37+00:00          Tur avsluttes før endeholdeplass    severe   KOL:Line:8\_1104  KOL:ServiceJourney:1104\_251007084239511\_3070  2026-10-24T00:00:00+02:00  2026-10-25T02:31:00+02:00
    \end{Verbatim}

    
    I check whether alert \texttt{journey\_ref} values correspond to GTFS
\texttt{trip\_id} values and whether \texttt{line\_ref} values
correspond to GTFS \texttt{route\_id} values. This confirms whether the
alert records can be linked to the trips and routes used in the
scheduled-versus-real-time comparison.

    \begin{Verbatim}[commandchars=\\\{\}]
alert journey\_ref found in gtfs trip\_id : 4/4
alert line\_ref found in gtfs route\_id   : 9/9
    \end{Verbatim}

    \subsection{\texorpdfstring{4. Real-time \texttt{stop\_visits}: SIRI
Estimated Timetable
data}{4. Real-time stop\_visits: SIRI Estimated Timetable data}}\label{real-time-stop_visits-siri-estimated-timetable-data}

I now examine the real-time source that records stop-level journey
information, including aimed, predicted, and actual times.

    \subsubsection{4.1 Loading the analysis
fields}\label{loading-the-analysis-fields}

The \texttt{stop\_visits} file contains approximately 10 million rows
and 26 columns. To reduce memory use, I read only the fields needed for
profiling, cleaning, and comparison. I do not rely on
\texttt{stop\_name} or \texttt{destination} here because they are mostly
empty in this extract; I use the GTFS stop table for stop names. I also
omit \texttt{source\_file} and \texttt{recorded\_at}, which are not
needed for the analysis below.

I interpret the main SIRI fields as follows:

\begin{itemize}
\tightlist
\item
  \texttt{data\_frame\_ref} identifies the journey's service date.
\item
  \texttt{journey\_ref}, \texttt{line\_ref}, and \texttt{stop\_ref} are
  checked against GTFS trip, route, and stop identifiers.
\item
  \texttt{aimed\_*} fields represent timetable times,
  \texttt{expected\_*} fields represent current predictions, and
  \texttt{actual\_*} fields represent recorded arrival or departure
  times.
\item
  \texttt{snapshot\_time} records when this version of the real-time row
  was captured.
\end{itemize}

    \begin{Verbatim}[commandchars=\\\{\}]
9,989,935 raw rows | 4,513 MB in memory
    \end{Verbatim}

            \begin{tcolorbox}[breakable, size=fbox, boxrule=.5pt, pad at break*=1mm, opacityfill=0]
\begin{Verbatim}[commandchars=\\\{\}]
              snapshot\_time data\_frame\_ref
journey\_ref          line\_ref direction\_ref vehicle\_mode monitored
journey\_cancelled call\_cancelled  call\_type  \textbackslash{}
0 2026-10-24 00:00:15+00:00     2026-10-24
KOL:ServiceJourney:1130\_250903083242106\_4040  KOL:Line:30\_1130       inbound
ferry      True             False          False  estimated
1 2026-10-24 00:00:15+00:00     2026-10-24
KOL:ServiceJourney:1130\_250903083242106\_4040  KOL:Line:30\_1130       inbound
ferry      True             False          False  estimated
2 2026-10-24 00:00:15+00:00     2026-10-24
KOL:ServiceJourney:5958\_250903083240300\_8028  KOL:Line:26\_5958       inbound
ferry      True             False          False  estimated

   stop\_order        stop\_ref              aimed\_arrival
expected\_arrival actual\_arrival            aimed\_departure
expected\_departure actual\_departure arrival\_status  \textbackslash{}
0           1  NSR:Quay:99269                        NaN
NaN            NaN  2026-10-24T18:30:00+02:00  2026-10-24T18:30:00+02:00
NaN            NaN
1           2  NSR:Quay:99704  2026-10-24T18:54:00+02:00
2026-10-24T18:54:00+02:00            NaN                        NaN
NaN              NaN            NaN
2           1  NSR:Quay:47191                        NaN
NaN            NaN  2026-10-24T16:40:00+02:00  2026-10-24T16:40:00+02:00
NaN            NaN

  departure\_status prediction\_inaccurate
0              NaN                 False
1              NaN                 False
2              NaN                 False
\end{Verbatim}
\end{tcolorbox}
        
    \subsubsection{Initial real-time data
assessment}\label{initial-real-time-data-assessment}

I focus on \texttt{stop\_visits} because it is the source used to
determine whether scheduled events have corresponding real-time
evidence. I first examine its identifiers, repeated records, and field
completeness before defining the cleaning and observation rules.

    \begin{Verbatim}[commandchars=\\\{\}]
Non-ferry records: 9,494,374  (raw file had 9,989,935)
  of which vehicle\_mode is NULL: 2,429
vehicle\_mode
bus    9491945
NaN       2429
Memory usage: 4,297 MB
    \end{Verbatim}

    
    \begin{Verbatim}[commandchars=\\\{\}]
                        True  False  Missing
journey\_cancelled       1686      0  9492688
call\_cancelled         64612      0  9429762
prediction\_inaccurate      0      0  9494374
    \end{Verbatim}

    
    
    \begin{Verbatim}[commandchars=\\\{\}]
                    Has value  Missing
aimed\_arrival         9231791   262583
expected\_arrival      6589616  2904758
actual\_arrival        2642175  6852199
aimed\_departure       9231791   262583
expected\_departure    6500328  2994046
actual\_departure      2731463  6762911
arrival\_status        6732697  2761677
departure\_status      6498192  2996182
    \end{Verbatim}

    
    \subsubsection{4.2 Understanding repeated
snapshots}\label{understanding-repeated-snapshots}

The initial count of distinct \texttt{stop\_name} values does not
represent the number of distinct stops because the feed's
\texttt{stop\_name} field is mostly empty. I count distinct
\texttt{stop\_ref} identifiers instead. I also measure how many rows are
recorded for each service-date, journey, and stop-order key. This shows
why raw row counts overstate the number of visits: the same visit can be
present in several snapshots.

    \begin{Verbatim}[commandchars=\\\{\}]
service dates in feed: ['2026-10-23', '2026-10-24', '2026-10-25', '2026-10-26',
'2026-10-27']
distinct stop\_ref  : 2,247

unique visits (journey × stop × date): 91,273
raw rows                              : 9,989,935
snapshots per visit:
count    91273.0
mean       109.5
std         67.2
min          2.0
25\%         93.0
50\%        105.0
75\%        125.0
max       1428.0
    \end{Verbatim}

    \begin{Verbatim}[commandchars=\\\{\}]
service dates in feed: ['2026-10-24', '2026-10-25', '2026-10-26', '2026-10-27']
distinct stop\_ref  : 2,225

unique visits (journey × stop × date): 90,261
raw buses                              : 9,494,374
snapshots per visit:
count    90261.0
mean       105.2
std         39.9
min          1.0
25\%         93.0
50\%        105.0
75\%        125.0
max       1428.0
    \end{Verbatim}

    \subsubsection{4.3 Assessing real-time data
quality}\label{assessing-real-time-data-quality}

I measure missing values and inspect the values in fields that describe
vehicle mode, call type, cancellation, monitoring, prediction quality,
and arrival or departure status. These checks help me understand which
fields are reliable enough to use in the cleaning rules and in the
interpretation of a visit.

    \begin{Verbatim}[commandchars=\\\{\}]
Share of nulls per column (\%):
journey\_cancelled        95.0
prediction\_inaccurate    95.0
call\_cancelled           94.4
actual\_arrival           73.4
actual\_departure         72.5
departure\_status         35.0
arrival\_status           32.6
expected\_departure       31.6
expected\_arrival         30.8
aimed\_departure           4.3
aimed\_arrival             4.3
snapshot\_time             0.0
data\_frame\_ref            0.0
direction\_ref             0.0
line\_ref                  0.0
journey\_ref               0.0
stop\_ref                  0.0
monitored                 0.0
stop\_order                0.0
call\_type                 0.0
vehicle\_mode              0.0

--- vehicle\_mode ---
vehicle\_mode
bus      9491945
ferry     495561
NaN         2429

--- call\_type ---
call\_type
estimated    7207575
recorded     2782360

--- arrival\_status ---
arrival\_status
delayed      6246114
NaN          3257238
early         420931
cancelled      64420
onTime          1232

--- departure\_status ---
departure\_status
delayed      6146452
NaN          3491743
onTime        290508
cancelled      61232

--- monitored ---
monitored
True     9989533
None         384
False         18

--- journey\_cancelled ---
journey\_cancelled
None     9492688
False     495463
True        1784

--- call\_cancelled ---
call\_cancelled
None     9429762
False     490055
True       70118

--- prediction\_inaccurate ---
prediction\_inaccurate
None     9494374
False     495561
    \end{Verbatim}

    \begin{Verbatim}[commandchars=\\\{\}]
Share of nulls per column (\%):
journey\_cancelled        100.0
prediction\_inaccurate    100.0
call\_cancelled            99.3
actual\_arrival            72.2
actual\_departure          71.2
departure\_status          31.6
expected\_departure        31.5
expected\_arrival          30.6
arrival\_status            29.1
aimed\_departure            2.8
aimed\_arrival              2.8
snapshot\_time              0.0
data\_frame\_ref             0.0
direction\_ref              0.0
line\_ref                   0.0
journey\_ref                0.0
stop\_ref                   0.0
monitored                  0.0
stop\_order                 0.0
call\_type                  0.0
vehicle\_mode               0.0

--- vehicle\_mode ---
vehicle\_mode
bus    9491945
NaN       2429

--- call\_type ---
call\_type
estimated    6735045
recorded     2759329

--- arrival\_status ---
arrival\_status
delayed      6246114
NaN          2761677
early         420931
cancelled      64420
onTime          1232

--- departure\_status ---
departure\_status
delayed      6146452
NaN          2996182
onTime        290508
cancelled      61232

--- monitored ---
monitored
True    9493990
None        384

--- journey\_cancelled ---
journey\_cancelled
None    9492688
True       1686

--- call\_cancelled ---
call\_cancelled
None    9429762
True      64612

--- prediction\_inaccurate ---
prediction\_inaccurate
None    9494374
    \end{Verbatim}

    \textbf{Basis for classifying a visit.} I use the presence of
\texttt{actual\_arrival} or \texttt{actual\_departure} as evidence that
the stop event was observed. I check how complete those fields are
before relying on this rule. If actual times are sparse, I report that
limitation and investigate whether another feed field provides a
defensible alternative rather than silently treating predictions as
completed visits.

            \begin{tcolorbox}[breakable, size=fbox, boxrule=.5pt, pad at break*=1mm, opacityfill=0]
\begin{Verbatim}[commandchars=\\\{\}]
0
\end{Verbatim}
\end{tcolorbox}
        
    \subsection{5. Reconstructing expected service from the
timetable}\label{reconstructing-expected-service-from-the-timetable}

I determine whether each GTFS service is active on each date by
combining its regular calendar with date-specific exceptions. A service
is active when the date falls within its calendar range and the weekday
flag is set. I then add services listed in \texttt{calendar\_dates} with
\texttt{exception\_type\ =\ 1} and remove those listed with
\texttt{exception\_type\ =\ 2}, such as holiday exceptions. Each
associated row in \texttt{stop\_times} represents one scheduled stop
event for an active trip.

    \begin{Verbatim}[commandchars=\\\{\}]
services per weekday flag (should NOT be all 0): \{'monday': 141, 'saturday': 16,
'sunday': 5\}
    \end{Verbatim}

    \subsubsection{5.1 Counting expected stop events by service
date}\label{counting-expected-stop-events-by-service-date}

I apply the calendar rules to identify the trips scheduled for each
service date, then count their stop-time records. These counts provide
the timetable-based baseline for comparison with real-time activity.

    
    \begin{Verbatim}[commandchars=\\\{\}]
   service\_date             route\_type\_name  trips  scheduled\_stop\_events  distinct\_stops
0    2026-10-24  High-speed passenger ferry     52                    391              44
1    2026-10-24                   Local bus   2537                  83983            2213
2    2026-10-24             Local car ferry    309                    722              28
3    2026-10-24       Local passenger ferry      2                     20              10
4    2026-10-24                   Night bus     65                   2775             864
5    2026-10-25  High-speed passenger ferry     53                    413              55
6    2026-10-25                   Local bus   1626                  51747            2038
7    2026-10-25             Local car ferry    315                    722              28
8    2026-10-25       Local passenger ferry      2                     20              10
9    2026-10-25                   Night bus     80                   3338             918
10   2026-10-25                  School bus      1                     10              10
11   2026-10-26  High-speed passenger ferry     97                    606              53
12   2026-10-26                   Local bus   4882                 147146            3799
13   2026-10-26             Local car ferry    371                    852              29
14   2026-10-26       Local passenger ferry      4                     40              10
15   2026-10-26                  School bus    511                   7457            3018
16   2026-10-27  High-speed passenger ferry     95                    601              53
17   2026-10-27                   Local bus   4898                 147702            3795
18   2026-10-27             Local car ferry    368                    865              40
19   2026-10-27       Local passenger ferry      4                     40              10
20   2026-10-27                  School bus    500                   7297            3015
    \end{Verbatim}

    
    \begin{Verbatim}[commandchars=\\\{\}]
BUS ONLY:
    \end{Verbatim}

    
    \begin{Verbatim}[commandchars=\\\{\}]
              trips  scheduled\_stop\_events  distinct\_stops
service\_date                                              
2026-10-24     2602                  86758            2247
2026-10-25     1707                  55095            2190
2026-10-26     5393                 154603            4673
2026-10-27     5398                 154999            4670
    \end{Verbatim}

    
    \begin{Verbatim}[commandchars=\\\{\}]
                   rows  trips            first\_snapshot
last\_snapshot
data\_frame\_ref
2026-10-23         2109      2 2026-10-24 00:00:15+00:00 2026-10-24
03:20:55+00:00
2026-10-24      9400601   2778 2026-10-24 00:00:15+00:00 2026-10-24
23:22:31+00:00
2026-10-25       523079    284 2026-10-24 00:00:15+00:00 2026-10-24
23:59:41+00:00
2026-10-26         5924      2 2026-10-24 00:00:15+00:00 2026-10-24
23:59:41+00:00
2026-10-27        58222      1 2026-10-24 01:14:34+00:00 2026-10-24
23:59:41+00:00
    \end{Verbatim}

    \begin{Verbatim}[commandchars=\\\{\}]
feed trips on 24th:         2778
  of those in GTFS trips:   2778
  of those I scheduled:     2778
2026-10-24 active service\_ids: 15 | trips on them: 2987
2026-10-26 active service\_ids: 121 | trips on them: 5883
service\_ids: trips 284 | calendar 223 | calendar\_dates 259
    \end{Verbatim}

    \begin{Verbatim}[commandchars=\\\{\}]
route\_type\_name
High-speed passenger ferry    36
Local bus                     75
Local car ferry               60
Local passenger ferry          2
Night bus                     14
Name: trip\_id, dtype: int64
    \end{Verbatim}

    \begin{Verbatim}[commandchars=\\\{\}]
departure\_time
08     1
09     2
10     1
11     1
15     2
17     1
18     1
20     1
21     1
24    66
25    12
Name: count, dtype: int64
    \end{Verbatim}

    \begin{Verbatim}[commandchars=\\\{\}]
0 missing trips use 0 service\_ids
    \end{Verbatim}

    
    \begin{Verbatim}[commandchars=\\\{\}]
Empty DataFrame
Columns: []
Index: []
    \end{Verbatim}

    
    \begin{Verbatim}[commandchars=\\\{\}]
monday       141
tuesday      128
wednesday    134
thursday     132
friday       144
saturday      16
sunday         5
dtype: int64
exception\_type
2    1942
1     780
Name: count, dtype: int64
date
20261023     6
20261024     5
20261025    12
20261026     9
20261027    11
dtype: int64
    \end{Verbatim}

    \begin{Verbatim}[commandchars=\\\{\}]
first departure hour (GTFS time):
Series([], )

active on 23rd / 25th:
Series([], )

scheduled-but-missing trips found in feed under another date:
data\_frame\_ref
2026-10-25    4
2026-10-26    1

first departure hour of the missing scheduled trips:
first\_dep
0      4
1      1
6      2
7      3
8      8
9      7
10     8
11     5
12     5
13     6
14     5
15     9
16     7
17     8
18     6
19     5
20     5
21     5
22     4
23     1
24    70
25    13
    \end{Verbatim}

    \texttt{expected\_per\_stop} summarizes the number of scheduled bus
visits at each stop on each service date. I use it as the
expected-activity baseline when assessing observed visits; it describes
the timetable, not the number of buses confirmed by the real-time feed.

    
    \begin{Verbatim}[commandchars=\\\{\}]
      service\_date         stop\_id  expected\_visits                stop\_name   stop\_lat  stop\_lon
11748   2026-10-27  NSR:Quay:46957              815             Vaktapoteket  58.966459  5.731226
7075    2026-10-26  NSR:Quay:46957              811             Vaktapoteket  58.966459  5.731226
11753   2026-10-27  NSR:Quay:46964              801             Vaktapoteket  58.966693  5.731406
7080    2026-10-26  NSR:Quay:46964              800             Vaktapoteket  58.966693  5.731406
1093    2026-10-24  NSR:Quay:46957              539             Vaktapoteket  58.966459  5.731226
1098    2026-10-24  NSR:Quay:46964              533             Vaktapoteket  58.966693  5.731406
11014   2026-10-27  NSR:Quay:46034              529            Vålandsbakken  58.957290  5.740290
6341    2026-10-26  NSR:Quay:46034              529            Vålandsbakken  58.957290  5.740290
7045    2026-10-26  NSR:Quay:46919              529              Statens hus  58.961118  5.737835
7106    2026-10-26  NSR:Quay:47005              529  Stavanger politistasjon  58.963432  5.735529
11779   2026-10-27  NSR:Quay:47005              529  Stavanger politistasjon  58.963432  5.735529
11718   2026-10-27  NSR:Quay:46919              529              Statens hus  58.961118  5.737835
11658   2026-10-27  NSR:Quay:46836              525           Lindahlsbakken  58.955635  5.741340
11013   2026-10-27  NSR:Quay:46033              525            Vålandsbakken  58.957532  5.739911
11777   2026-10-27  NSR:Quay:47003              525  Stavanger politistasjon  58.962876  5.735773
    \end{Verbatim}

    
    \subsection{6. Preparing and classifying real-time
visits}\label{preparing-and-classifying-real-time-visits}

I clean the real-time records so that each usable row represents one
journey-stop event rather than one downloaded snapshot. I record row
removals by rule in \texttt{removal\_log}, which lets me report how each
cleaning decision affects the data.

My cleaning process is:

\begin{enumerate}
\def\labelenumi{\arabic{enumi}.}
\tightlist
\item
  Keep the selected service dates and remove rows without the
  identifiers needed to define a visit.
\item
  Consolidate repeated snapshots using the latest snapshot for each
  visit. Because actual arrival and departure times may have been
  recorded in an earlier snapshot, I retain those values when the latest
  snapshot does not contain them. Other fields, including predictions
  and cancellation flags, come from the latest snapshot.
\item
  Parse timestamps as UTC and derive local-time fields only for
  interpretation.
\item
  Check whether trip, route, and stop identifiers link to the static
  GTFS tables.
\item
  Exclude ferry services and retain bus route types and stops that exist
  in GTFS.
\item
  Classify each visit, calculate delay from a single selected arrival or
  departure event, and add analysis flags and time fields.
\end{enumerate}

    \begin{Verbatim}[commandchars=\\\{\}]
keep chosen service dates                                  9,494,374 ->
9,494,374   removed            0
drop rows with null key / stop\_ref                         9,494,374 ->
9,494,374   removed            0
visits whose newest snapshot has no actual time but an older one does: 43
(rescued below)
de-duplicate snapshots (keep latest per visit)             9,494,374 ->
90,261   removed    9,404,113
    \end{Verbatim}

    \begin{Verbatim}[commandchars=\\\{\}]
OK   aimed\_arrival: values that failed to parse: 0 of 87,571
OK   expected\_arrival: values that failed to parse: 0 of 2,828
OK   actual\_arrival: values that failed to parse: 0 of 84,786
OK   aimed\_departure: values that failed to parse: 0 of 87,571
OK   expected\_departure: values that failed to parse: 0 of 2,726
OK   actual\_departure: values that failed to parse: 0 of 84,888
aimed\_arrival         97.0
expected\_arrival       3.1
actual\_arrival        93.9
aimed\_departure       97.0
expected\_departure     3.0
actual\_departure      94.0
timestamp dtype: datetime64[us, UTC]
    \end{Verbatim}

    \begin{Verbatim}[commandchars=\\\{\}]
journey\_ref in GTFS trips : 100.0\%
stop\_ref in GTFS stops    : 100.0\%
line\_ref in GTFS routes   : 100.0\%
OK   duplicate route\_id in routes: 0 of 318
OK   duplicate stop\_id in stops: 0 of 7,676
ferry lines: 5 by vehicle\_mode, 21 by route\_type, 21 together
exclude ferry lines (vehicle\_mode or route\_type)              90,261 ->
90,092   removed          169

route\_type\_name  Local bus  Night bus  School bus
vehicle\_mode
bus                  85098       4990           4
keep bus route types only (3, 700-799)                        90,092 ->
90,092   removed            0
keep stops that exist in GTFS                                 90,092 ->
90,092   removed            0
    \end{Verbatim}

    \subsubsection{6.4 Defining an observed stop
visit}\label{defining-an-observed-stop-visit}

I classify each cleaned visit using the following rules:

\begin{longtable}[]{@{}
  >{\raggedright\arraybackslash}p{(\linewidth - 2\tabcolsep) * \real{0.5000}}
  >{\raggedright\arraybackslash}p{(\linewidth - 2\tabcolsep) * \real{0.5000}}@{}}
\toprule\noalign{}
\begin{minipage}[b]{\linewidth}\raggedright
\texttt{visit\_status}
\end{minipage} & \begin{minipage}[b]{\linewidth}\raggedright
Rule
\end{minipage} \\
\midrule\noalign{}
\endhead
\bottomrule\noalign{}
\endlastfoot
\texttt{cancelled} & The journey or stop call is marked as cancelled. \\
\texttt{observed} & It is not cancelled and an actual arrival or
departure time is present. \\
\texttt{unconfirmed} & It is not cancelled, but no actual time is
present. The visit may be in the future or the feed may never report an
actual time. \\
\end{longtable}

I include only \texttt{observed} visits in the confirmed-visit table,
while retaining unconfirmed visits to show where real-time records exist
without an actual time.

For each visit, I use one event consistently for \texttt{visit\_time},
\texttt{aimed\_time}, and \texttt{delay\_s}: departure when an actual
departure is available, otherwise arrival. This accommodates first stops
with no arrival and last stops with no departure. \texttt{delay\_event}
identifies which event was used.

I retain the following flags rather than removing those records:
\texttt{delay\_outlier} marks delays more than three hours late or more
than 30 minutes early; \texttt{on\_time} marks delays from 60 seconds
early to 180 seconds late and is meaningful only for observed visits;
\texttt{no\_prediction} marks a call without an expected time; and
\texttt{is\_actual} records whether \texttt{call\_type} is
\texttt{recorded}, allowing comparison with the actual-time
classification.

I store timestamps in UTC. The \texttt{local\_date},
\texttt{local\_hour}, \texttt{weekday}, and \texttt{visit\_time\_local}
fields are derived in \texttt{Europe/Oslo}. On 2026-10-25, the local
02:00 hour occurs twice because of the daylight-saving transition, so I
use \texttt{hour\_utc} for unambiguous hourly counts across that date. I
use \texttt{data\_frame\_ref} for the service date because a journey
shortly after midnight may belong to the previous service day.

    \begin{Verbatim}[commandchars=\\\{\}]
is\_actual (call\_type == recorded) vs has actual\_* time:
has\_actual\_time  False  True
is\_actual
False             2784     43
True                 0  87265
    \end{Verbatim}

    
    \begin{Verbatim}[commandchars=\\\{\}]
              visits  share
visit\_status               
observed       87279  0.969
unconfirmed     2432  0.027
cancelled        381  0.004
    \end{Verbatim}

    
    \begin{Verbatim}[commandchars=\\\{\}]
delay outliers among observed: 21
on time among observed with a delay: 62.0\%
    \end{Verbatim}

    
    \begin{Verbatim}[commandchars=\\\{\}]
                                               rule  rows\_before  rows\_after  removed
0                         keep chosen service dates      9494374     9494374        0
1                drop rows with null key / stop\_ref      9494374     9494374        0
2    de-duplicate snapshots (keep latest per visit)      9494374       90261  9404113
3  exclude ferry lines (vehicle\_mode or route\_type)        90261       90092      169
4            keep bus route types only (3, 700-799)        90092       90092        0
5                     keep stops that exist in GTFS        90092       90092        0
    \end{Verbatim}

    
    \subsection{7. Comparing scheduled and real-time stop
events}\label{comparing-scheduled-and-real-time-stop-events}

I compare the scheduled bus events with cleaned real-time records to
assess which scheduled stops are observed, cancelled, unconfirmed, or
absent from the feed. I summarize the outcomes by stop, route, and
service date.

    I match each scheduled stop event to a real-time visit using service
date, trip, stop, and the stop occurrence within that trip. I use the
occurrence number rather than matching \texttt{stop\_sequence} directly
to \texttt{stop\_order}, because those numbering fields are not
guaranteed to use the same sequence. Before interpreting the comparison,
I check how many scheduled and real-time trip identifiers overlap.

    \begin{Verbatim}[commandchars=\\\{\}]
scheduled bus trips : 15,100
trips seen in ET    : 2,674
in both             : 2,674
ET only (not in static timetable, e.g. added trips): 0
timetable only (no real-time data at all)         : 12,426
    \end{Verbatim}

    
    \begin{Verbatim}[commandchars=\\\{\}]
final\_status  cancelled  observed  trip\_not\_in\_feed  unconfirmed  scheduled\_total  observed\_share
service\_date                                                                                     
2026-10-24          324     83724              2534          176            86758           0.965
2026-10-25           17      3548             49274         2256            55095           0.064
2026-10-26            3         1            154599            0           154603           0.000
2026-10-27           37         6            154956            0           154999           0.000
    \end{Verbatim}

    
    \textbf{Interpreting the comparison.} The outcome categories describe
the evidence available in the two sources:

\begin{itemize}
\tightlist
\item
  \texttt{observed}: a scheduled event has a matched real-time visit
  with an actual time.
\item
  \texttt{cancelled}: the matched real-time record indicates
  cancellation.
\item
  \texttt{unconfirmed}: a matching real-time record exists, but it has
  no actual time.
\item
  \texttt{trip\_not\_in\_feed}: the timetable includes the trip, but the
  real-time feed contains no record for it. This may reflect limited
  feed coverage, not necessarily a missed bus.
\item
  \texttt{stop\_not\_in\_feed}: the trip is present in the real-time
  feed, but the scheduled stop has no corresponding real-time record.
\end{itemize}

I interpret these categories separately so that feed coverage gaps are
not confused with confirmed cancellations or failures to operate.

    
    \begin{Verbatim}[commandchars=\\\{\}]
     service\_date          stop\_id  expected\_visits  observed\_visits  cancelled  median\_delay\_s                stop\_name   stop\_lat  stop\_lon  observed\_share
1093   2026-10-24   NSR:Quay:46957              539              519          3      103.299500             Vaktapoteket  58.966459  5.731226           0.963
1098   2026-10-24   NSR:Quay:46964              533              517          3      152.329000             Vaktapoteket  58.966693  5.731406           0.970
1070   2026-10-24   NSR:Quay:46920              343              329          1       58.331000              Statens hus  58.960533  5.738111           0.959
1118   2026-10-24   NSR:Quay:47003              343              329          1       72.383000  Stavanger politistasjon  58.962876  5.735773           0.959
675    2026-10-24   NSR:Quay:46033              343              329          1       64.317000            Vålandsbakken  58.957532  5.739911           0.959
1027   2026-10-24   NSR:Quay:46836              343              329          1       51.329000           Lindahlsbakken  58.955635  5.741340           0.959
676    2026-10-24   NSR:Quay:46034              338              328          2       75.314500            Vålandsbakken  58.957290  5.740290           0.970
1069   2026-10-24   NSR:Quay:46919              338              328          2       79.307000              Statens hus  58.961118  5.737835           0.970
1120   2026-10-24   NSR:Quay:47005              338              328          2       57.826500  Stavanger politistasjon  58.963432  5.735529           0.970
1231   2026-10-24   NSR:Quay:47271              326              313          1       40.355000               Strømsbrua  58.953664  5.741909           0.960
1230   2026-10-24   NSR:Quay:47270              322              312          2       87.861500               Strømsbrua  58.954434  5.742140           0.969
1106   2026-10-24   NSR:Quay:46974              308              293          1       21.776031            Jernbaneveien  58.966973  5.733415           0.951
22     2026-10-24  NSR:Quay:101964              270              259          2      139.332000       St. Svithuns plass  58.962902  5.725116           0.959
1279   2026-10-24   NSR:Quay:47383              271              259          1      109.321000                   Kannik  58.964585  5.728295           0.956
604    2026-10-24   NSR:Quay:45855              270              259          2      122.326000            Holbergs gate  58.961471  5.716522           0.959
    \end{Verbatim}

    
    
    \begin{Verbatim}[commandchars=\\\{\}]
    service\_date route\_short\_name  expected\_visits  observed\_visits  observed\_share
189   2026-10-26                5             9106                0           0.000
466   2026-10-27                5             9106                0           0.000
475   2026-10-27                6             8342                0           0.000
198   2026-10-26                6             8342                0           0.000
40    2026-10-24                5             7974             7810           0.979
452   2026-10-27                3             7304                0           0.000
175   2026-10-26                3             7274                0           0.000
459   2026-10-27                4             7008                0           0.000
44    2026-10-24                6             6976             6820           0.978
182   2026-10-26                4             6880                0           0.000
399   2026-10-27                1             6858                0           0.000
122   2026-10-26                1             6792                0           0.000
35    2026-10-24                4             6362             5946           0.935
461   2026-10-27               42             6266                0           0.000
184   2026-10-26               42             6171                0           0.000
    \end{Verbatim}

    
    \begin{center}
    \adjustimage{max size={0.9\linewidth}{0.9\paperheight}}{building_silver_layer_v2_files/building_silver_layer_v2_59_1.png}
    \end{center}
    { \hspace*{\fill} \\}
    
    \subsection{8. Checking if the BUS actually stopped at the stop or
not}\label{checking-if-the-bus-actually-stopped-at-the-stop-or-not}

    \begin{Verbatim}[commandchars=\\\{\}]
observed visits: 87,279 | with a dwell (both times): 82,119
negative dwell: 0

count    82119.0
mean         0.0
std          0.0
min          0.0
5\%           0.0
25\%          0.0
50\%          0.0
75\%          0.0
95\%          0.0
max          0.0

visits per dwell second (0-20 s):
dwell\_s
0.0    82119
    \end{Verbatim}

    \begin{center}
    \adjustimage{max size={0.9\linewidth}{0.9\paperheight}}{building_silver_layer_v2_files/building_silver_layer_v2_61_1.png}
    \end{center}
    { \hspace*{\fill} \\}
    
    \begin{Verbatim}[commandchars=\\\{\}]
               visits  share
stop\_class
likely\_passed   82119  0.941
unknown          5160  0.059

share passed by local hour:
            visits  share\_passed
local\_hour
0             2671           1.0
1             2465           1.0
2             1453           1.0
3             1063           1.0
4              542           1.0
5              349           1.0
6              765           1.0
7             1867           1.0
8             2664           1.0
9             3363           1.0
10            4195           1.0
11            5254           1.0
12            5249           1.0
13            5372           1.0
14            5370           1.0
15            5378           1.0
16            5041           1.0
17            4989           1.0
18            4818           1.0
19            4357           1.0
20            3818           1.0
21            3759           1.0
22            3628           1.0
23            3689           1.0
    \end{Verbatim}

    
    \begin{Verbatim}[commandchars=\\\{\}]
           stop\_ref            stop\_name  visits  share\_passed
3   NSR:Quay:100428  Coop Extra Moksheim      33           1.0
4   NSR:Quay:100429  Coop Extra Moksheim      34           1.0
5   NSR:Quay:100430        Skarhaugvegen      34           1.0
12  NSR:Quay:101344            Resahagen      56           1.0
16  NSR:Quay:101964   St. Svithuns plass     271           1.0
24  NSR:Quay:102676        Kveldsbeleika      31           1.0
25  NSR:Quay:102677        Kveldsbeleika      30           1.0
30  NSR:Quay:102744             Endrestø      74           1.0
51  NSR:Quay:105300              Iglamyr      49           1.0
63  NSR:Quay:107636       Nådlandsbråtet     203           1.0
69  NSR:Quay:107871         Løwenstrasse     189           1.0
81  NSR:Quay:108770           Helsehuset     111           1.0
82  NSR:Quay:108874       SUS Ullandhaug     157           1.0
83  NSR:Quay:108875       SUS Ullandhaug     157           1.0
84  NSR:Quay:108881                Jåtun     106           1.0
    \end{Verbatim}

    
    \begin{Verbatim}[commandchars=\\\{\}]
call\_cancelled visits: 307

arrival\_status:
 arrival\_status
NaN          83800
cancelled      307
delayed         93
early           24

departure\_status:
 departure\_status
NaN          83865
cancelled      291
delayed         68
    \end{Verbatim}

    \subsection{8. Saving the analysis
tables}\label{saving-the-analysis-tables}

I save the cleaned visit data and scheduled-versus-real-time comparison
so that the next stage can reuse the results without repeating the full
preparation process:

\begin{itemize}
\tightlist
\item
  \texttt{clean\_stop\_visits.parquet} contains one row per confirmed
  bus visit, with stop, time, delay, and route information. Timestamps
  are in UTC, with derived \texttt{Europe/Oslo} fields and a DST-safe
  \texttt{hour\_utc} field.
\item
  \texttt{stop\_event\_status.parquet} contains each scheduled bus stop
  event and its comparison outcome.
\item
  \texttt{expected\_per\_stop.parquet} contains the timetable-based
  expected number of visits by stop and service date.
\item
  \texttt{removal\_log.csv} records the rows removed by each cleaning
  rule, providing an audit trail for the report.
\end{itemize}

    \begin{Verbatim}[commandchars=\\\{\}]
clean\_stop\_visits: 87,279 rows
stored timestamp dtype: datetime64[ns, UTC]
    \end{Verbatim}

    
    \begin{Verbatim}[commandchars=\\\{\}]
  service\_date                                   journey\_ref         line\_ref route\_short\_name route\_type\_name direction\_ref  stop\_order        stop\_ref               stop\_name   stop\_lat  stop\_lon  \textbackslash{}
0   2026-10-24  KOL:ServiceJourney:1109\_251007078405577\_1140  KOL:Line:8\_1109                9       Local bus             0           1  NSR:Quay:46967  Stavanger bussterminal  58.966548  5.733496   
1   2026-10-24  KOL:ServiceJourney:1109\_251007078405577\_1140  KOL:Line:8\_1109                9       Local bus             0           2  NSR:Quay:46955               Klubbgata  58.970015  5.734989   
2   2026-10-24  KOL:ServiceJourney:1109\_251007078405577\_1140  KOL:Line:8\_1109                9       Local bus             0           3  NSR:Quay:46413             Jorenholmen  58.971445  5.737626   
3   2026-10-24  KOL:ServiceJourney:1109\_251007078405577\_1140  KOL:Line:8\_1109                9       Local bus             0           4  NSR:Quay:48551              Fiskepiren  58.971420  5.741196   
4   2026-10-24  KOL:ServiceJourney:1109\_251007078405577\_1140  KOL:Line:8\_1109                9       Local bus             0           5  NSR:Quay:48605             Bekhuskaien  58.971015  5.744256   

              aimed\_arrival                actual\_arrival           aimed\_departure              actual\_departure                    visit\_time                aimed\_time delay\_event  delay\_s  \textbackslash{}
0                       NaT                           NaT 2026-10-23 23:30:00+00:00 2026-10-23 23:31:49.296+00:00 2026-10-23 23:31:49.296+00:00 2026-10-23 23:30:00+00:00   departure  109.296   
1 2026-10-23 23:33:00+00:00 2026-10-23 23:33:30.302+00:00 2026-10-23 23:33:00+00:00 2026-10-23 23:33:30.302+00:00 2026-10-23 23:33:30.302+00:00 2026-10-23 23:33:00+00:00   departure   30.302   
2 2026-10-23 23:33:00+00:00 2026-10-23 23:34:54.327+00:00 2026-10-23 23:33:00+00:00 2026-10-23 23:34:54.327+00:00 2026-10-23 23:34:54.327+00:00 2026-10-23 23:33:00+00:00   departure  114.327   
3 2026-10-23 23:34:00+00:00 2026-10-23 23:35:39.327+00:00 2026-10-23 23:34:00+00:00 2026-10-23 23:35:39.327+00:00 2026-10-23 23:35:39.327+00:00 2026-10-23 23:34:00+00:00   departure   99.327   
4 2026-10-23 23:34:00+00:00 2026-10-23 23:36:05.316+00:00 2026-10-23 23:34:00+00:00 2026-10-23 23:36:05.316+00:00 2026-10-23 23:36:05.316+00:00 2026-10-23 23:34:00+00:00   departure  125.316   

   delay\_arrival\_s  delay\_departure\_s  delay\_outlier  on\_time  is\_actual  n\_snapshots  local\_date  local\_hour   weekday                  hour\_utc              visit\_time\_local  
0              NaN            109.296          False     True       True            2  2026-10-24           1  Saturday 2026-10-23 23:00:00+00:00 2026-10-24 01:31:49.296+02:00  
1           30.302             30.302          False     True       True            2  2026-10-24           1  Saturday 2026-10-23 23:00:00+00:00 2026-10-24 01:33:30.302+02:00  
2          114.327            114.327          False     True       True            2  2026-10-24           1  Saturday 2026-10-23 23:00:00+00:00 2026-10-24 01:34:54.327+02:00  
3           99.327             99.327          False     True       True            2  2026-10-24           1  Saturday 2026-10-23 23:00:00+00:00 2026-10-24 01:35:39.327+02:00  
4          125.316            125.316          False     True       True            2  2026-10-24           1  Saturday 2026-10-23 23:00:00+00:00 2026-10-24 01:36:05.316+02:00  
    \end{Verbatim}

    
    \subsection{9. Report-ready
visualizations}\label{report-ready-visualizations}

The figures below summarize scheduled service, real-time feed coverage,
observed delays, and stop behavior. They use the analysis tables built
above and are saved as high-resolution PNG files in
\texttt{clean/report\_figures/} for use in a written report. A low
observed share or a record missing from the feed indicates limited
evidence in these data; it does not, by itself, prove that a bus did not
run.
 
            
    
    \begin{center}
    \adjustimage{max size={0.9\linewidth}{0.9\paperheight}}{building_silver_layer_v2_files/building_silver_layer_v2_67_0.png}
    \end{center}
    { \hspace*{\fill} \\}
    

    \begin{center}
    \adjustimage{max size={0.9\linewidth}{0.9\paperheight}}{building_silver_layer_v2_files/building_silver_layer_v2_67_1.png}
    \end{center}
    { \hspace*{\fill} \\}
    
    \begin{Verbatim}[commandchars=\\\{\}, frame=single, framerule=2mm, rulecolor=\color{outerrorbackground}]
\textcolor{ansi-red}{---------------------------------------------------------------------------}
\textcolor{ansi-red}{KeyError}                                  Traceback (most recent call last)
\textcolor{ansi-cyan}{Cell}\textcolor{ansi-cyan}{ }\textcolor{ansi-green}{In[41]}\textcolor{ansi-green}{, line 52}
\textcolor{ansi-green}{     48} axes[\textcolor{ansi-green}{1}, \textcolor{ansi-green}{0}].set\_ylabel(\textcolor{ansi-yellow}{"Observed visits"})
\textcolor{ansi-green}{     49} 
\textcolor{ansi-green}{     50} \def\tcRGB{\textcolor[RGB]}\expandafter\tcRGB\expandafter{\detokenize{95,135,135}}{\# The short-dwell threshold is the PASS\_MAX\_S definition from section 8.2.}
\textcolor{ansi-green}{     51} known\_visits = v[v[\textcolor{ansi-yellow}{"stop\_class"}].ne(\textcolor{ansi-yellow}{"unknown"})]
\textcolor{ansi-green}{---> }\textcolor{ansi-green}{52} pass\_by\_hour = known\_visits.groupby(\textcolor{ansi-yellow}{"local\_hour"})[\textcolor{ansi-yellow}{"passed"}].agg(visits=\textcolor{ansi-yellow}{"size"}, share\_passed=\textcolor{ansi-yellow}{"mean"})
\textcolor{ansi-green}{     53} axes[\textcolor{ansi-green}{1}, \textcolor{ansi-green}{1}].bar(pass\_by\_hour.index, pass\_by\_hour[\textcolor{ansi-yellow}{"visits"}], color=\textcolor{ansi-yellow}{"\#cbd5e1"}, label=\textcolor{ansi-yellow}{"Classified visits"})
\textcolor{ansi-green}{     54} pass\_ax = axes[\textcolor{ansi-green}{1}, \textcolor{ansi-green}{1}].twinx()
\textcolor{ansi-green}{     55} pass\_ax.plot(pass\_by\_hour.index, pass\_by\_hour["share\_passed"] * 100,

\textcolor{ansi-cyan}{File }\textcolor{ansi-green}{\textasciitilde{}\textbackslash{}AppData\textbackslash{}Local\textbackslash{}Python\textbackslash{}pythoncore-3.14-64\textbackslash{}Lib\textbackslash{}site-packages\textbackslash{}pandas\textbackslash{}core\textbackslash{}groupby\textbackslash{}generic.py:3002}, in \textcolor{ansi-cyan}{DataFrameGroupBy.\_\_getitem\_\_}\textcolor{ansi-blue}{(self, key)}
\textcolor{ansi-green}{   2995} \def\tcRGB{\textcolor[RGB]}\expandafter\tcRGB\expandafter{\detokenize{0,135,0}}{\textbf{if}} \def\tcRGB{\textcolor[RGB]}\expandafter\tcRGB\expandafter{\detokenize{0,135,0}}{isinstance}(key, \def\tcRGB{\textcolor[RGB]}\expandafter\tcRGB\expandafter{\detokenize{0,135,0}}{tuple}) \def\tcRGB{\textcolor[RGB]}\expandafter\tcRGB\expandafter{\detokenize{175,0,255}}{\textbf{and}} \def\tcRGB{\textcolor[RGB]}\expandafter\tcRGB\expandafter{\detokenize{0,135,0}}{len}(key) > \textcolor{ansi-green}{1}:
\textcolor{ansi-green}{   2996}     \def\tcRGB{\textcolor[RGB]}\expandafter\tcRGB\expandafter{\detokenize{95,135,135}}{\# if len == 1, then it becomes a SeriesGroupBy and this is actually}
\textcolor{ansi-green}{   2997}     \def\tcRGB{\textcolor[RGB]}\expandafter\tcRGB\expandafter{\detokenize{95,135,135}}{\# valid syntax, so don't raise}
\textcolor{ansi-green}{   2998}     \def\tcRGB{\textcolor[RGB]}\expandafter\tcRGB\expandafter{\detokenize{0,135,0}}{\textbf{raise}} \def\tcRGB{\textcolor[RGB]}\expandafter\tcRGB\expandafter{\detokenize{215,95,95}}{\textbf{ValueError}}(
\textcolor{ansi-green}{   2999}         \textcolor{ansi-yellow}{"}\textcolor{ansi-yellow}{Cannot subset columns with a tuple with more than one element. }\textcolor{ansi-yellow}{"}
\textcolor{ansi-green}{   3000}         \textcolor{ansi-yellow}{"}\textcolor{ansi-yellow}{Use a list instead.}\textcolor{ansi-yellow}{"}
\textcolor{ansi-green}{   3001}     )
\textcolor{ansi-green}{-> }\textcolor{ansi-green}{3002} \def\tcRGB{\textcolor[RGB]}\expandafter\tcRGB\expandafter{\detokenize{0,135,0}}{\textbf{return}} \textcolor{ansi-black}{\setlength{\fboxsep}{0pt}\colorbox{ansi-yellow}{super\strut}}\textcolor{ansi-black}{\setlength{\fboxsep}{0pt}\colorbox{ansi-yellow}{(\strut}}\textcolor{ansi-black}{\setlength{\fboxsep}{0pt}\colorbox{ansi-yellow}{)\strut}}\textcolor{ansi-black}{\setlength{\fboxsep}{0pt}\colorbox{ansi-yellow}{.\strut}}\textcolor{ansi-black}{\setlength{\fboxsep}{0pt}\colorbox{ansi-yellow}{\_\_getitem\_\_\strut}}\textcolor{ansi-black}{\setlength{\fboxsep}{0pt}\colorbox{ansi-yellow}{(\strut}}\textcolor{ansi-black}{\setlength{\fboxsep}{0pt}\colorbox{ansi-yellow}{key\strut}}\textcolor{ansi-black}{\setlength{\fboxsep}{0pt}\colorbox{ansi-yellow}{)\strut}}

\textcolor{ansi-cyan}{File }\textcolor{ansi-green}{\textasciitilde{}\textbackslash{}AppData\textbackslash{}Local\textbackslash{}Python\textbackslash{}pythoncore-3.14-64\textbackslash{}Lib\textbackslash{}site-packages\textbackslash{}pandas\textbackslash{}core\textbackslash{}base.py:241}, in \textcolor{ansi-cyan}{SelectionMixin.\_\_getitem\_\_}\textcolor{ansi-blue}{(self, key)}
\textcolor{ansi-green}{    239} \def\tcRGB{\textcolor[RGB]}\expandafter\tcRGB\expandafter{\detokenize{0,135,0}}{\textbf{else}}:
\textcolor{ansi-green}{    240}     \def\tcRGB{\textcolor[RGB]}\expandafter\tcRGB\expandafter{\detokenize{0,135,0}}{\textbf{if}} key \def\tcRGB{\textcolor[RGB]}\expandafter\tcRGB\expandafter{\detokenize{175,0,255}}{\textbf{not}} \def\tcRGB{\textcolor[RGB]}\expandafter\tcRGB\expandafter{\detokenize{175,0,255}}{\textbf{in}} \def\tcRGB{\textcolor[RGB]}\expandafter\tcRGB\expandafter{\detokenize{0,135,0}}{self}.obj:
\textcolor{ansi-green}{--> }\textcolor{ansi-green}{241}         \def\tcRGB{\textcolor[RGB]}\expandafter\tcRGB\expandafter{\detokenize{0,135,0}}{\textbf{raise}} \def\tcRGB{\textcolor[RGB]}\expandafter\tcRGB\expandafter{\detokenize{215,95,95}}{\textbf{KeyError}}(\textcolor{ansi-yellow}{f}\textcolor{ansi-yellow}{"}\textcolor{ansi-yellow}{Column not found: }\def\tcRGB{\textcolor[RGB]}\expandafter\tcRGB\expandafter{\detokenize{175,95,135}}{\textbf{\{}}key\def\tcRGB{\textcolor[RGB]}\expandafter\tcRGB\expandafter{\detokenize{175,95,135}}{\textbf{\}}}\textcolor{ansi-yellow}{"})
\textcolor{ansi-green}{    242}     ndim = \def\tcRGB{\textcolor[RGB]}\expandafter\tcRGB\expandafter{\detokenize{0,135,0}}{self}.obj[key].ndim
\textcolor{ansi-green}{    243}     \def\tcRGB{\textcolor[RGB]}\expandafter\tcRGB\expandafter{\detokenize{0,135,0}}{\textbf{return}} \def\tcRGB{\textcolor[RGB]}\expandafter\tcRGB\expandafter{\detokenize{0,135,0}}{self}.\_gotitem(key, ndim=ndim)

\textcolor{ansi-red}{KeyError}: 'Column not found: passed'
    \end{Verbatim}

    \begin{center}
    \adjustimage{max size={0.9\linewidth}{0.9\paperheight}}{building_silver_layer_v2_files/building_silver_layer_v2_68_1.png}
    \end{center}
    { \hspace*{\fill} \\}
    
    \subsection{What I can do from here}\label{what-i-can-do-from-here}

The current work establishes the timetable baseline, cleans the
available real-time records, and compares the two sources. I can build
on it in the following ways:

\begin{enumerate}
\def\labelenumi{\arabic{enumi}.}
\tightlist
\item
  \textbf{Measure feed coverage.} Calculate the proportion of scheduled
  trips and stop events represented in the real-time feed by service
  date, route, and operator. Use coverage-aware denominators when
  reporting observed shares so that missing feed data is not mistaken
  for poor service.
\item
  \textbf{Describe bus punctuality.} For confirmed visits, summarize
  delay distributions and on-time shares by route, stop, time of day,
  and service date. Report medians and ranges alongside averages because
  extreme delays can distort the mean.
\item
  \textbf{Investigate cancellations and missing events.} Compare
  cancellation outcomes and feed gaps across routes and dates, and check
  whether patterns are concentrated in particular parts of the network.
\item
  \textbf{Add service-alert context.} Relate alerts to the affected
  routes or trips and examine whether they coincide with cancellations
  or unusual delays. Treat this as descriptive evidence rather than
  proof that an alert caused an outcome.
\item
  \textbf{Present the findings.} Section 9 adds report-ready figures for
  service volume, feed coverage, punctuality, stop behavior, and data
  quality. The report should state the study dates, observation rules,
  and the number of events excluded or unconfirmed.
\end{enumerate}

Before drawing operational conclusions, I should confirm that the chosen
dates and routes have adequate real-time coverage and review the delay
outliers and the distinction between recorded actual times and other
feed indicators.


    % Add a bibliography block to the postdoc
    
    
    
\end{document}
